\documentclass[10pt,a4paper]{amsart}
\usepackage{amsmath,amssymb,amsthm,graphicx}
% Retained minimal preamble for the Stacks Project chapter varieties.tex.
% Verbatim excerpt of the chapter preamble: Xy-pic setup (native preamble lines 14-18), the native theorem
% declarations with the mathematical macros (native preamble lines 154-217), and the project own no-op
% definitions of the reference/slogan/history display environments (native wrapper.tex lines 1-3) that the
% counted statement uses for its slogan. Copied mechanically; nothing added, removed or redefined.
% For commutative diagrams we use Xy-pic
\usepackage[all]{xy}

% We use 2cell for 2-commutative diagrams.
\xyoption{2cell}
% Theorem environments.
%
\theoremstyle{plain}
\newtheorem{theorem}[subsection]{Theorem}
\newtheorem{proposition}[subsection]{Proposition}
\newtheorem{lemma}[subsection]{Lemma}

\theoremstyle{definition}
\newtheorem{definition}[subsection]{Definition}
\newtheorem{example}[subsection]{Example}
\newtheorem{exercise}[subsection]{Exercise}
\newtheorem{situation}[subsection]{Situation}

\theoremstyle{remark}
\newtheorem{remark}[subsection]{Remark}
\newtheorem{remarks}[subsection]{Remarks}

\numberwithin{equation}{subsection}

% Macros
%
\def\lim{\mathop{\mathrm{lim}}\nolimits}
\def\colim{\mathop{\mathrm{colim}}\nolimits}
\def\Spec{\mathop{\mathrm{Spec}}}
\def\Hom{\mathop{\mathrm{Hom}}\nolimits}
\def\Ext{\mathop{\mathrm{Ext}}\nolimits}
\def\SheafHom{\mathop{\mathcal{H}\!\mathit{om}}\nolimits}
\def\SheafExt{\mathop{\mathcal{E}\!\mathit{xt}}\nolimits}
\def\Sch{\mathit{Sch}}
\def\Mor{\mathop{\mathrm{Mor}}\nolimits}
\def\Ob{\mathop{\mathrm{Ob}}\nolimits}
\def\Sh{\mathop{\mathit{Sh}}\nolimits}
\def\NL{\mathop{N\!L}\nolimits}
\def\CH{\mathop{\mathrm{CH}}\nolimits}
\def\proetale{{pro\text{-}\acute{e}tale}}
\def\etale{{\acute{e}tale}}
\def\QCoh{\mathit{QCoh}}
\def\Ker{\mathop{\mathrm{Ker}}}
\def\Im{\mathop{\mathrm{Im}}}
\def\Coker{\mathop{\mathrm{Coker}}}
\def\Coim{\mathop{\mathrm{Coim}}}

% Boxtimes
%
\DeclareMathSymbol{\boxtimes}{\mathbin}{AMSa}{"02}

%
% Macros for moduli stacks/spaces
%
\def\QCohstack{\mathcal{QC}\!\mathit{oh}}
\def\Cohstack{\mathcal{C}\!\mathit{oh}}
\def\Spacesstack{\mathcal{S}\!\mathit{paces}}
\def\Quotfunctor{\mathrm{Quot}}
\def\Hilbfunctor{\mathrm{Hilb}}
\def\Curvesstack{\mathcal{C}\!\mathit{urves}}
\def\Polarizedstack{\mathcal{P}\!\mathit{olarized}}
\def\Complexesstack{\mathcal{C}\!\mathit{omplexes}}
% \Pic is the operator that assigns to X its picard group, usage \Pic(X)
% \Picardstack_{X/B} denotes the Picard stack of X over B
% \Picardfunctor_{X/B} denotes the Picard functor of X over B
\def\Pic{\mathop{\mathrm{Pic}}\nolimits}
\def\Picardstack{\mathcal{P}\!\mathit{ic}}
\def\Picardfunctor{\mathrm{Pic}}
\def\Deformationcategory{\mathcal{D}\!\mathit{ef}}
\newenvironment{reference}{}{}
\newenvironment{slogan}{}{}
\newenvironment{history}{}{}
\usepackage{cleveref}
\crefname{originallemma}{original lemma}{original lemmas}
\Crefname{originallemma}{Original lemma}{Original lemmas}
\makeatletter
\newlabel{algebra-lemma-separably-closed-connected-implies-geometric}{{037U}{source}{}{sourcepointer-1}{}}
\newlabel{algebra-lemma-separably-closed-connected-implies-geometric@cref}{{[originallemma][0][]037U}{[1][0][]source}{}{}{}}
\makeatother
\crefname{originallemma}{original lemma}{original lemmas}
\Crefname{originallemma}{Original lemma}{Original lemmas}
\makeatletter
\newlabel{lemma-affine-geometrically-connected}{{0386}{source}{}{sourcepointer-2}{}}
\newlabel{lemma-affine-geometrically-connected@cref}{{[originallemma][0][]0386}{[1][0][]source}{}{}{}}
\makeatother
\crefname{originallemma}{original lemma}{original lemmas}
\Crefname{originallemma}{Original lemma}{Original lemmas}
\makeatletter
\newlabel{lemma-bijection-connected-components}{{0385}{source}{}{sourcepointer-3}{}}
\newlabel{lemma-bijection-connected-components@cref}{{[originallemma][0][]0385}{[1][0][]source}{}{}{}}
\makeatother
\title{Base change from a separably algebraically closed field induces a bijection on connected components}
\author{The Stacks Project Authors}
\date{stacks-project:varieties.tex:local-snapshot-20261011; GNU FDL 1.2 or later}
\begin{document}
\maketitle
\noindent\textit{Editorial note.} All delivered mathematical text is the author's own, copied from the pinned Stacks Project chapter file varieties.tex; no statement or proof step was rewritten, reordered or removed.\par\medskip
\noindent\textit{Editorial note.} The preamble is the chapter's own retained minimal preamble (sources/stacks-varieties/retained-preamble.tex), a verbatim excerpt of the native Xy-pic setup, of the native theorem declarations with mathematical macros, and of the project's own no-op definitions of the reference display environments; no macro was invented or redefined.\par\medskip
\noindent\textit{Editorial note.} The two same-chapter references are located at their \textbackslash{}label lines inside the pinned chapter file; the cross-chapter pointer is located in the project's own tag list staged in the snapshot (sources/stacks-chow/tags.txt), at the verbatim line carrying its canonical tag.\par\medskip
\noindent\textit{Allowed standard prerequisites (editorial).}\begin{itemize}
\item Over a separably algebraically closed field, a field extension is geometrically connected. Algebra, Lemma labelled algebra-lemma-separably-closed-connected-implies-geometric (canonical Stacks tag 037U); invoked at proof lines 781-784. Supplies that k' is geometrically connected over k; the passage to the spectrum and the product decomposition are retained in the counted span.
\item The spectrum of a geometrically connected algebra is geometrically connected. Same chapter, Lemma labelled lemma-affine-geometrically-connected (source line 758, Stacks tag 0386); invoked at proof lines 785-786. Supplies that Z = Spec(k') is geometrically connected over k, so that the component-bijection lemma applies.
\item Projection from a geometrically connected factor induces a bijection on connected components. Same chapter, Lemma labelled lemma-bijection-connected-components (source line 732, Stacks tag 0385, counted separately in this batch); invoked at proof lines 787-788. Supplies the bijection on components once the factor is geometrically connected; recognizing the base change as a product with Z is the counted step.
\end{itemize}\medskip
\ifdefined\hypertarget\hypertarget{sourcepointer-1}{}\fi
\noindent\textit{Source reference.} Original lemma 037U: sources/stacks-chow/tags.txt, source lines 4213--4213. Cross-chapter standard prerequisite invoked once at proof lines 781-784 for geometric connectedness of k' over k; the scheme-theoretic assembly on top of it is retained in the counted span.\par\medskip
\ifdefined\hypertarget\hypertarget{sourcepointer-2}{}\fi
\noindent\textit{Source reference.} Original lemma 0386: sources/stacks-varieties/pinned-original-source.tex, source lines 758--758. Same-chapter pointer invoked at proof lines 785-786 to convert geometric connectedness of k' into geometric connectedness of Spec(k').\par\medskip
\ifdefined\hypertarget\hypertarget{sourcepointer-3}{}\fi
\noindent\textit{Source reference.} Original lemma 0385: sources/stacks-varieties/pinned-original-source.tex, source lines 732--732. Same-chapter pointer invoked at proof lines 787-788 to obtain the bijection on components; the identification X\_\{k'\} = Z x\_k X is retained in the counted span.\par\medskip
\begin{lemma}
\label{lemma-separably-closed-field-connected-components}
Let $k'/k$ be an extension of fields.
Let $X$ be a scheme over $k$.
Assume $k$ separably algebraically closed.
Then the morphism $X_{k'} \to X$ induces a bijection of connected
components. In particular, $X$ is geometrically connected over $k$
if and only if $X$ is connected.
\end{lemma}

\begin{proof}
Since $k$ is separably algebraically closed we see that
$k'$ is geometrically connected over $k$, see
Algebra,
Lemma \ref{algebra-lemma-separably-closed-connected-implies-geometric}.
Hence $Z = \Spec(k')$ is geometrically connected over $k$ by
Lemma \ref{lemma-affine-geometrically-connected}
above. Since $X_{k'} = Z \times_k X$ the result is a special case of
Lemma \ref{lemma-bijection-connected-components}.
\end{proof}

\end{document}
